% ─────────────────────────────────────────────────────────────────────────────
%  Appendices — each \chapter{} becomes Appendix A, B, C, etc.
% ─────────────────────────────────────────────────────────────────────────────

\chapter{Repository Structure}
\label{app:repo}

The SHOOTRZ repository is organised as follows:

\begin{lstlisting}[language={},caption={Top-level repository structure.},label={lst:repo}]
SHOOTRZ/
|-- backend/
|   |-- main.py                  # App factory, lifespan, routers
|   |-- config/mvp_config.yaml   # Tunable pipeline parameters
|   |-- routers/                 # 9 HTTP router modules
|   |-- mvp/core/                # 6-stage pipeline
|   |   |-- pipeline.py          # Orchestrator
|   |   |-- video_loader.py      # Frame extraction + stride
|   |   |-- pose_estimation.py   # MediaPipe wrapper
|   |   |-- signal_smoothing.py  # Savitzky-Golay + interpolation
|   |   |-- shot_detection.py    # Multi-signal state machine
|   |   |-- angle_computation.py # Joint angle calculation
|   |   `-- metrics.py           # Normative scoring
|   |-- inference/               # Phase detector, ball tracker
|   |-- metrics/normative_ranges.json  # Research-backed thresholds
|   |-- services/
|   |   |-- mvp_job_service.py   # Job orchestration
|   |   |-- job_store.py         # SQLite job state (72 h TTL)
|   |   `-- llm/                 # Gemini integration
|   |-- recommender/             # FAISS + LinUCB
|   |-- feedback/                # Rule-based coaching cues
|   |-- chat/context_builder.py  # LLM context injection
|   |-- storage/db.py            # Supabase ORM layer
|   |-- contracts/               # Pydantic request/response models
|   |-- tests/                   # pytest test suite
|   `-- requirements.txt         # Python dependencies
|
|-- src/                         # React Native application
|   |-- screens/                 # 11 application screens
|   |-- components/              # 30+ UI components
|   |-- services/                # API, Supabase, chat, storage clients
|   |-- context/                 # Auth, History, Profile contexts
|   |-- types/contracts.ts       # Canonical TypeScript interfaces
|   |-- theme/                   # Design tokens, typography, motion
|   `-- navigation/AppNavigator.tsx
|
|-- supabase/
|   |-- schema_complete.sql      # Canonical schema (13 tables)
|   `-- migration_coach_context_rpc.sql
|
|-- App.tsx                      # Root component
|-- package.json                 # Frontend dependencies
`-- README.md
\end{lstlisting}

% ─────────────────────────────────────────────────────────────────────────────
\chapter{Key API Endpoints}
\label{app:api}

Table~\ref{tab:api_endpoints} documents the key API endpoints exposed by the SHOOTRZ backend.

\begin{table}[H]
  \centering
  \caption{SHOOTRZ backend API endpoints.}
  \label{tab:api_endpoints}
  \rowcolors{2}{gray!8}{white}
  \begin{tabularx}{\textwidth}{p{1.5cm}p{6.4cm}p{1cm}X}
    \toprule
    \textbf{Method} & \textbf{Endpoint} & \textbf{Auth} & \textbf{Description} \\
    \midrule
    POST   & \texttt{/mvp/analyze}                       & No           & Upload video; receive \texttt{job\_id} \\
    GET    & \texttt{/mvp/result/\{job\_id\}}            & No           & Poll analysis status and result \\
    GET    & \texttt{/mvp/artifact/\{run\_id\}/\{file\}} & No           & Download overlay video, CSVs \\
    POST   & \texttt{/api/analysis/complete}             & \textbf{Yes} & Persist analysis result to Supabase \\
    POST   & \texttt{/api/chat}                          & \textbf{Yes} & SSE streaming coaching chat \\
    GET    & \texttt{/api/user/analysis-history}         & No*          & Paginated analysis history \\
    DELETE & \texttt{/api/user/account}                  & No*          & Full account deletion \\
    GET    & \texttt{/api/recommend}                     & No*          & Personalised drill recommendation \\
    GET    & \texttt{/health}                            & No           & Version, uptime, Gemini status \\
    \bottomrule
  \end{tabularx}
\end{table}

\noindent{\small * Routes marked No* lack authentication middleware; remediation is documented in Section~\ref{sec:future_work}.}

% ─────────────────────────────────────────────────────────────────────────────
\chapter{Configuration Parameters}
\label{app:config}

Table~\ref{tab:config} lists the tunable parameters in \lstinline|backend/config/mvp_config.yaml| with their production defaults.

\begin{table}[H]
  \centering
  \caption{Key \texttt{mvp\_config.yaml} parameters.}
  \label{tab:config}
  \rowcolors{2}{gray!8}{white}
  \begin{tabularx}{\textwidth}{ll>{\raggedright\arraybackslash}X}
    \toprule
    \textbf{Parameter} & \textbf{Default} & \textbf{Description} \\
    \midrule
    \lstinline|pose_detection.model_complexity| & 1 & MediaPipe model size (0 fastest, 2 most accurate) \\
    \lstinline|pose_detection.min_detection_confidence| & 0.5 & Minimum per-landmark confidence \\
    \lstinline|pose_detection.confidence_threshold| & 0.3 & Frame-level emission threshold \\
    \lstinline|video.max_frames| & 300 & Hard cap on frames processed (stride-adjusted) \\
    \lstinline|video.min_duration| & 1.0\,s & Shortest acceptable clip \\
    \lstinline|smoothing.window_length| & 5 & Savitzky--Golay window (must be odd) \\
    \lstinline|smoothing.polyorder| & 2 & Savitzky--Golay polynomial order \\
    \lstinline|shot_detection.knee_flexion_threshold| & 100.0\textdegree & Adaptive crouch detection seed threshold \\
    \lstinline|shot_detection.use_legacy_detector| & false & Fall back to simple heuristic \\
    \lstinline|output.save_overlay_video| & true & Generate annotated skeleton video \\
    \bottomrule
  \end{tabularx}
\end{table}

% ─────────────────────────────────────────────────────────────────────────────
\chapter{Normative Range Reference}
\label{app:normative}

Table~\ref{tab:normative_full} lists all normative ranges used for metric scoring, sourced from peer-reviewed sports biomechanics literature.

\begin{table}[H]
  \centering
  \caption{Normative ranges for primary biomechanical metrics.}
  \label{tab:normative_full}
  \rowcolors{2}{gray!8}{white}
  \begin{tabularx}{\textwidth}{lccX}
    \toprule
    \textbf{Metric} & \textbf{Target Range} & \textbf{Optimal Range} & \textbf{Source} \\
    \midrule
    Elbow extension at release & 165--180\textdegree & 170--180\textdegree & \cite{Cabarkapa2021} \\
    Elbow flexion (preparatory) & 70--85\textdegree & 75--85\textdegree & \cite{Cabarkapa2021} \\
    Knee flexion at crouch & 100--120\textdegree & 105--115\textdegree & \cite{Cabarkapa2021} \\
    Forearm verticality & 0--10\textdegree & 0--8\textdegree & \cite{Cabarkapa2021} \\
    Wrist follow-through & 15--30\textdegree & 18--28\textdegree & \cite{Struzik2014} \\
    Release angle at 2.8\,m & 74--84\textdegree & 76--82\textdegree & \cite{Okazaki2012} \\
    Release angle at 4.6\,m & 60--70\textdegree & 63--68\textdegree & \cite{Okazaki2012} \\
    Release angle at 6.0\,m+ & 55--65\textdegree & 58--62\textdegree & \cite{Okazaki2012} \\
    Ball entry angle & 45--55\textdegree & 48--52\textdegree & --- \\
    \bottomrule
  \end{tabularx}
\end{table}

% ─────────────────────────────────────────────────────────────────────────────
\chapter{Scoring Weight Distribution}
\label{app:weights}

Table~\ref{tab:weights} lists the metric weights used in the confidence-weighted geometric mean overall score computation.

\begin{table}[H]
  \centering
  \caption{Metric weights in the overall score computation.}
  \label{tab:weights}
  \rowcolors{2}{gray!8}{white}
  \begin{tabular}{lcc}
    \toprule
    \textbf{Metric} & \textbf{Weight} & \textbf{Normative Key} \\
    \midrule
    Elbow extension & 0.32 & \lstinline|elbow_flexion_release| \\
    Knee bend & 0.24 & \lstinline|knee_flexion| \\
    Release angle & 0.20 & \lstinline|release_angle_4.6m| \\
    Wrist follow-through & 0.14 & \lstinline|wrist_follow_through| \\
    Forearm verticality & 0.10 & \lstinline|forearm_verticality| \\
    \bottomrule
  \end{tabular}
\end{table}

\noindent Weights are re-normalised per analysis over the subset of metrics with a valid selected frame. When ball tracking is disabled, the release angle weight is excluded and remaining weights are rescaled to sum to 1.0.

% ─────────────────────────────────────────────────────────────────────────────
\chapter{Environment Variables Reference}
\label{app:env}

\subsection*{Backend (\texttt{SHOOTRZ/backend/.env})}

\begin{table}[H]
  \centering
  \caption{Backend environment variables.}
  \label{tab:backend_env}
  \rowcolors{2}{gray!8}{white}
  \begin{tabularx}{\textwidth}{llX}
    \toprule
    \textbf{Variable} & \textbf{Default} & \textbf{Description} \\
    \midrule
    \lstinline|SUPABASE_URL| & --- & Supabase project URL \\
    \lstinline|SUPABASE_SERVICE_KEY| & --- & Service-role key (bypasses RLS) \\
    \lstinline|GEMINI_API_KEY| & --- & Google GenAI API key \\
    \lstinline|GEMINI_MODEL| & \lstinline|gemini-2.5-flash| & LLM model identifier \\
    \lstinline|SHOOTRZ_MAX_CONCURRENT| & 8 & Analysis job semaphore cap \\
    \lstinline|SHOOTRZ_ENABLE_BALL| & 0 & Enable YOLOv8n ball tracking \\
    \lstinline|SHOOTRZ_MAX_UPLOAD_MB| & 200 & Upload size cap \\
    \lstinline|SHOOTRZ_AI_TIMEOUT_S| & 10 & Gemini call timeout (seconds) \\
    \lstinline|SHOOTRZ_SAVE_OVERLAY| & (from YAML) & Override overlay video generation \\
    \bottomrule
  \end{tabularx}
\end{table}

\subsection*{Frontend (\texttt{SHOOTRZ/.env})}

\begin{table}[H]
  \centering
  \caption{Frontend environment variables.}
  \label{tab:frontend_env}
  \rowcolors{2}{gray!8}{white}
  \begin{tabularx}{\textwidth}{lX}
    \toprule
    \textbf{Variable} & \textbf{Description} \\
    \midrule
    \lstinline|EXPO_PUBLIC_SUPABASE_URL| & Supabase project URL \\
    \lstinline|EXPO_PUBLIC_SUPABASE_ANON_KEY| & Supabase anonymous (public) key \\
    \lstinline|EXPO_PUBLIC_API_URL| & Backend base URL (e.g., \lstinline|http://192.168.1.x:8000|) \\
    \bottomrule
  \end{tabularx}
\end{table}
